\chapter{Mouvement brownien}\label{ch:qm:brownian-motion}

Un trader achète une action dont la volatilité annuelle est de 30~\%, et place un stop-loss 2~\% sous le
cours d'entrée, puis demande au desk de risque quelle est la probabilité que le stop soit atteint dans le mois. La
probabilité que l'action \emph{termine} le mois simplement sous le stop vaut 41~\%~; celle qu'elle
touche le stop à un instant quelconque du mois est exactement le double, soit 82~\%. Si le stop est
contrôlé seulement à la clôture de chaque journée, la réponse est 72~\%, et une simulation qui observe les
clôtures et oublie ce qui se passe entre elles annonce 72~\% pour un contrat qui est en réalité
déclenché dans 82~\% des cas. Les trois nombres proviennent d'un même objet, le \omterm{def:qm:brownian-motion:bm}{mouvement brownien}, et
de trois de ses propriétés~: sa loi à un instant fixé, le principe de réflexion pour son
minimum courant, et le \omterm{def:qm:brownian-motion:bridge}{pont brownien} qui le décrit entre deux observations. Ce
chapitre construit l'objet, démontre les propriétés utilisées dans la série, et se termine par la
\omterm{def:qm:brownian-motion:qv}{variation quadratique}, la propriété qui fait du calcul stochastique un outil différent du
calcul ordinaire.

\section{Des marches aléatoires au mouvement brownien}

\begin{definition}[Mouvement brownien, processus gaussien]\label{def:qm:brownian-motion:bm}
Un \emph{processus gaussien}\index{processus gaussien} est un processus $(X_t)$ dont les lois de dimension finie
$(X_{t_1}, \dots, X_{t_k})$ sont toutes normales multivariées~; sa loi est fixée par sa moyenne
$m(t)$ et sa covariance $c(s,t)$. Un \emph{mouvement brownien}\index{mouvement brownien} standard
$(W_t)_{t \ge 0}$ est un \omterm{def:qm:probability-at-speed:filtration}{processus adapté} tel que $W_0 = 0$, à trajectoires continues, et dont les accroissements
$W_t - W_s$ sont indépendants de $\mathcal F_s$ et suivent $\mathcal N(0, t - s)$ pour $s < t$.
\end{definition}

De manière équivalente, $W$ est le \omterm{def:qm:brownian-motion:bm}{processus gaussien} centré à trajectoires continues et de covariance
$c(s,t) = \min(s,t)$~: pour $s < t$, $\Cov(W_s, W_t) = \Cov(W_s, W_s + (W_t - W_s)) = s$.
L'existence n'a rien d'évident, puisque la définition impose la continuité d'une infinité non dénombrable
de variables aléatoires à la fois. Wiener donna la première construction en 1923~; celle qui sert
à la simulation procède du grossier au fin (la \cref{sec:qm:brownian-motion:bridge}). La raison
pour laquelle le \omterm{def:qm:brownian-motion:bm}{mouvement brownien} apparaît partout tient à un théorème central limite pour les trajectoires.

\begin{theorem}[Principe d'invariance de Donsker]\label{thm:qm:brownian-motion:donsker}
Soient $\xi_1, \xi_2, \dots$ des variables indépendantes et identiquement distribuées, centrées et
de variance égale à un, $S_k = \sum_{i \le k}\xi_i$, et $X^{(n)}_t = n^{-1/2}(S_{\lfloor nt\rfloor} +
(nt - \lfloor nt\rfloor)\xi_{\lfloor nt\rfloor + 1})$ la marche interpolée linéairement et réduite.
Alors $X^{(n)} \xrightarrow{d} W$ comme éléments aléatoires de $C[0, 1]$~: $\E[F(X^{(n)})] \to
\E[F(W)]$ pour toute fonctionnelle continue bornée $F$ de la trajectoire.
\end{theorem}

\admitted

Le théorème de Donsker justifie l'usage du \omterm{def:qm:brownian-motion:bm}{mouvement brownien} pour des prix construits à partir de
nombreuses petites transactions indépendantes, quelle que soit leur loi individuelle~: le maximum de la
trajectoire, le temps passé sous un niveau, le prix moyen, toutes les fonctionnelles continues, convergent ensemble. Il
ne le justifie pas pour des prix gouvernés par quelques grands sauts (chapitre~6).

\begin{proposition}[Invariances]\label{prop:qm:brownian-motion:scaling}
Si $W$ est un \omterm{def:qm:brownian-motion:bm}{mouvement brownien}, il en va de même pour $-W_t$~; $c^{-1/2}W_{ct}$ pour tout $c > 0$
(\emph{changement d'échelle})~; $W_{s+t} - W_s$ pour $s$ fixé~; et $tW_{1/t}$ (valant 0 en $t = 0$).
\end{proposition}

\begin{proof}
Chacun est un \omterm{def:qm:brownian-motion:bm}{processus gaussien} centré de covariance $\min(s,t)$~; pour le changement d'échelle, $c^{-1}\min(cs,
ct) = \min(s,t)$~; pour l'inversion, $st\min(1/s,1/t) = \min(s,t)$. La continuité en 0 de $tW_{1/t}$
découle de la loi forte $W_u/u \to 0$.
\end{proof}

Le changement d'échelle est la règle de la racine carrée du temps qu'un trader applique sans y penser~: une
action à volatilité annuelle de 30~\% a un écart type sur un jour de $0.30/\sqrt{252} = 1.9\%$ et un
sur un mois, $0.30\sqrt{21/252} = 8.7\%$. La \cref{fig:qm:brownian-motion:paths} montre six
trajectoires à l'intérieur de l'enveloppe $\pm 2\sqrt t$.

\begin{omfigure}
\begin{tikzpicture}
\begin{axis}[width=11cm,height=5cm,
  xlabel={$t$ (années)},ylabel={$W_t$},
  tick label style={font=\small},label style={font=\small},
  xmin=0,xmax=1,ymin=-2.4,ymax=2.4,grid=major,grid style={gray!25},
  table/col sep=comma]
\addplot[gray,dashed,thick] table[x=t,y=up]{figdata/methods/02-brownian-motion/paths.csv};
\addplot[gray,dashed,thick] table[x=t,y=dn]{figdata/methods/02-brownian-motion/paths.csv};
\addplot[omDef,thin] table[x=t,y=w0]{figdata/methods/02-brownian-motion/paths.csv};
\addplot[omThm,thin] table[x=t,y=w1]{figdata/methods/02-brownian-motion/paths.csv};
\addplot[omProp,thin] table[x=t,y=w2]{figdata/methods/02-brownian-motion/paths.csv};
\addplot[omExo,thin] table[x=t,y=w3]{figdata/methods/02-brownian-motion/paths.csv};
\addplot[black!60,thin] table[x=t,y=w4]{figdata/methods/02-brownian-motion/paths.csv};
\addplot[omDef!50,thin] table[x=t,y=w5]{figdata/methods/02-brownian-motion/paths.csv};
\end{axis}
\end{tikzpicture}

\omcaption{Six trajectoires browniennes sur $[0,1]$ (250 pas) à l'intérieur de l'enveloppe $\pm 2\sqrt t$ (tirets), qui
contient 95~\% de la loi à chaque instant fixé. L'écart croît comme $\sqrt t$, et non comme $t$.
Données~: le tutoriel du chapitre, avec graine fixée.}\label{fig:qm:brownian-motion:paths}
\end{omfigure}

\section{Propriétés des trajectoires et propriété de Markov}

Les trajectoires browniennes sont continues et nulle part régulières. Presque sûrement, elles sont höldériennes
d'ordre quelconque $\alpha < \tfrac12$ et d'aucun ordre $\alpha > \tfrac12$, et dérivables en
aucun point~; elles traversent tout niveau qu'elles atteignent une infinité de fois juste après. Rien de
cela n'est une pathologie gratuite~: c'est la raison pour laquelle une couverture rééquilibrée en continu dans le temps
conserve un coût non nul (chapitre~3), et la raison pour laquelle une barrière touchée une fois est
touchée plusieurs fois.

\begin{definition}[Processus de Markov, propriété de Markov forte]\label{def:qm:brownian-motion:markov}
Un \omterm{def:qm:probability-at-speed:filtration}{processus adapté} $(X_t)$ est un \emph{processus de Markov}\index{processus de Markov} si pour tous $s, t
\ge 0$ et toute fonction mesurable bornée $f$, $\E[f(X_{s+t}) \mid \mathcal F_s] = \E[f(X_{s+t}) \mid X_s]$~:
connaissant le présent, le passé n'apporte rien. Il possède la \emph{propriété de Markov forte}\index{propriété de Markov forte}
si la même égalité vaut lorsque $s$ est remplacé par un \omterm{def:qm:probability-at-speed:stopping}{temps d'arrêt} fini quelconque $\tau$.
\end{definition}

\begin{theorem}[Le mouvement brownien est fortement markovien]\label{thm:qm:brownian-motion:strongmarkov}
Pour tout \omterm{def:qm:probability-at-speed:stopping}{temps d'arrêt} $\tau$ tel que $\P(\tau < \infty) = 1$, le processus $B_t = W_{\tau + t} -
W_\tau$ est un \omterm{def:qm:brownian-motion:bm}{mouvement brownien} indépendant de $\mathcal F_\tau$.
\end{theorem}

\admitted

La démonstration approche $\tau$ par au-dessus, au moyen de \omterm{def:qm:probability-at-speed:stopping}{temps d'arrêt} à valeurs dénombrables, pour
lesquels l'énoncé se ramène à la propriété de Markov simple appliquée à chaque valeur (Karatzas et
Shreve, 1991, §2.6). La \omterm{def:qm:brownian-motion:markov}{propriété de Markov forte} est ce qui permet à la section suivante de fonctionner~: après
la première fois qu'une trajectoire atteint un niveau, la suite est un \omterm{def:qm:brownian-motion:bm}{mouvement brownien} neuf, aussi
probable de monter que de descendre.

\section{Principe de réflexion et temps de premier passage}

\begin{definition}[Temps de premier passage]\label{def:qm:brownian-motion:passage}
Pour un niveau $b$, le \emph{temps de premier passage}\index{temps de premier passage} (ou \emph{temps
d'atteinte}\index{temps d'atteinte}) d'un processus continu $X$ est $\tau_b = \inf\{t \ge 0 : X_t = b\}$~;
c'est un \omterm{def:qm:probability-at-speed:stopping}{temps d'arrêt}.
\end{definition}

\begin{theorem}[Principe de réflexion]\label{thm:qm:brownian-motion:reflection}
Pour $b > 0$ et $t > 0$, posons $M_t = \max_{s \le t}W_s$. Alors
\[
\P(M_t \ge b) = \P(\tau_b \le t) = 2\P(W_t \ge b) = 2\bigl(1 - \Phi(b/\sqrt t)\bigr),
\]
et plus précisément $\P(M_t \ge b, W_t \le a) = \P(W_t \ge 2b - a)$ pour $a \le b$.
\end{theorem}

\begin{proof}
Sur $\{\tau_b \le t\}$, on définit la trajectoire réfléchie $\tilde W_s = W_s$ pour $s \le \tau_b$ et $\tilde
W_s = 2b - W_s$ ensuite (la \cref{fig:qm:brownian-motion:reflect}). Par
le \cref{thm:qm:brownian-motion:strongmarkov}, $W_{\tau_b + u} - b$ est un \omterm{def:qm:brownian-motion:bm}{mouvement brownien}
indépendant de $\mathcal F_{\tau_b}$, et il en va de même de son opposé~; donc $\tilde W$ est un
\omterm{def:qm:brownian-motion:bm}{mouvement brownien}. L'événement $\{M_t \ge b, W_t \le a\}$ pour $W$ est l'événement $\{\tilde W_t \ge 2b - a\}$ pour
$\tilde W$ (les trajectoires atteignent $b$ au même instant), qui a la probabilité de $\{W_t \ge 2b -
a\}$. Avec $a = b$~: $\P(M_t \ge b) = \P(M_t \ge b, W_t \le b) + \P(W_t > b) = 2\P(W_t \ge b)$.
\end{proof}

\begin{omfigure}
\begin{tikzpicture}
\begin{axis}[width=11cm,height=5cm,
  xlabel={$t$},ylabel={$W_t$},
  tick label style={font=\small},label style={font=\small},
  xmin=0,xmax=1,ymin=-0.6,ymax=1.3,grid=major,grid style={gray!25},
  legend columns=2,legend style={font=\footnotesize,at={(0.5,-0.3)},anchor=north,draw=none,fill=none,/tikz/every even column/.append style={column sep=8pt}},
  table/col sep=comma]
\addplot[omThm,thick,dashed] table[x=t,y=refl]{figdata/methods/02-brownian-motion/reflect.csv};
\addplot[omDef,thick] table[x=t,y=w]{figdata/methods/02-brownian-motion/reflect.csv};
\addplot[black,thin,forget plot] coordinates {(0,0.6) (1,0.6)};
\addplot[gray,thin,forget plot] coordinates {(0.618,-0.6) (0.618,1.3)};
\node[font=\footnotesize,anchor=south west,fill=white,inner sep=1pt] at (axis cs:0.01,0.63) {niveau $b = 0.6$};
\node[font=\footnotesize,anchor=north west,fill=white,inner sep=1pt] at (axis cs:0.64,-0.33) {$\tau_b = 0.618$};
\legend{réfléchie après $\tau_b$,trajectoire}
\end{axis}
\end{tikzpicture}

\omcaption{Le principe de réflexion. Après le premier passage en $b$, la trajectoire (trait plein) et son
image miroir par rapport à $b$ (tirets) sont des prolongements aussi probables, de sorte que toute trajectoire qui finit sous
$b$ après l'avoir touchée est associée à une trajectoire qui finit au-dessus. Données~: une trajectoire générée à partir d'une graine fixée du
tutoriel.}\label{fig:qm:brownian-motion:reflect}
\end{omfigure}

En dérivant par rapport à $t$, on obtient la loi du \omterm{def:qm:brownian-motion:passage}{temps de premier passage}~: $\tau_b$ admet pour densité
$b(2\pi t^3)^{-1/2}e^{-b^2/2t}$, finie presque sûrement mais d'espérance infinie, car la densité
décroît comme $t^{-3/2}$. Un prix sans dérive finit par atteindre tout niveau, et il ne faut pas
l'attendre. Le stop de l'exemple introductif est le principe de réflexion appliqué au logarithme du prix.

\begin{definition}[Mouvement brownien avec dérive]\label{def:qm:brownian-motion:drift}
Un \emph{mouvement brownien avec dérive}\index{mouvement brownien avec dérive} de dérive $\mu$ et de volatilité
$\sigma > 0$ est $X_t = \mu t + \sigma W_t$.
\end{definition}

\begin{proposition}[Premier passage avec dérive]\label{prop:qm:brownian-motion:driftpassage}
Pour $X_t = \mu t + \sigma W_t$ et $b < 0$,
\[
\P\bigl(\min_{s \le T}X_s \le b\bigr) = \Phi\Bigl(\frac{b - \mu T}{\sigma\sqrt T}\Bigr) +
e^{2\mu b/\sigma^2}\,\Phi\Bigl(\frac{b + \mu T}{\sigma\sqrt T}\Bigr).
\]
\end{proposition}

\admitted

La démonstration change de mesure pour supprimer la dérive et applique le principe de réflexion
(un exercice du chapitre~5 l'établit). Pour l'exemple introductif,
$b = \ln 0.98 = -0.0202$ et $\sigma\sqrt T = 0.0866$~: sans dérive, le stop est touché avec
probabilité $2\Phi(-0.233) = 81.6\%$~; avec la dérive $-\sigma^2/2$ qui fait du prix lui-même
une \omterm{def:qm:probability-at-speed:martingale}{martingale}, 82.4~\%.

\begin{proposition}[Suivi discret~: le décalage de Broadie--Glasserman--Kou]\label{prop:qm:brownian-motion:bgk}
Si $X$ n'est observé qu'aux instants $k\Delta t$, la probabilité qu'une observation soit
inférieure ou égale à $b < 0$ est, au premier ordre en $\sqrt{\Delta t}$, la probabilité continue pour le
niveau décalé $b - \beta\sigma\sqrt{\Delta t}$, avec $\beta = -\zeta(\tfrac12)/\sqrt{2\pi} \approx
0.5826$.
\end{proposition}

\admitted

Le décalage est le dépassement moyen de la trajectoire au-delà du niveau, à la première observation située
sous celui-ci (Broadie, Glasserman et Kou, 1997). Avec des clôtures quotidiennes, $\beta\sigma\sqrt{\Delta t} =
0.0110$~: le stop à 2~\% se comporte comme un stop à 3.1~\% suivi en continu, et il est touché avec
une probabilité de 71.9~\%, contre 72.0~\% dans une simulation de 200\,000 mois de clôtures quotidiennes.
La \Cref{fig:qm:brownian-motion:stops} présente les trois courbes pour des distances de stop allant de 1~\% à
10~\%. L'erreur du suivi discret ne décroît qu'en $\sqrt{\Delta t}$~: affiner une simulation quotidienne
pour la porter à quatre, seize et soixante-quatre observations par jour donne 77.1~\%, 79.1~\% et 80.5~\%
contre 81.6~\% en continu, l'écart se réduisant à peu près de moitié à chaque fois que le pas est divisé par
quatre. Une simulation qui veut la réponse continue ne doit pas affiner la grille~; elle doit utiliser
le bridge de la \cref{sec:qm:brownian-motion:bridge}.

\begin{omfigure}
\begin{tikzpicture}
\begin{axis}[width=11cm,height=5cm,
  xlabel={distance du stop sous le point d'entrée (\%)},ylabel={P(stop touché en 21 jours)},
  tick label style={font=\small},label style={font=\small},
  xmin=0.5,xmax=10.5,ymin=0,ymax=1,grid=major,grid style={gray!25},
  legend columns=2,legend style={font=\footnotesize,at={(0.5,-0.3)},anchor=north,draw=none,fill=none,/tikz/every even column/.append style={column sep=8pt}},
  table/col sep=comma]
\addplot[omDef,very thick] table[x=dist,y=cont]{figdata/methods/02-brownian-motion/stops.csv};
\addplot[omDef,only marks,mark=o,mark size=2.4pt] table[x=dist,y=bridge]{figdata/methods/02-brownian-motion/stops.csv};
\addplot[omThm,very thick,dashed] table[x=dist,y=bgk]{figdata/methods/02-brownian-motion/stops.csv};
\addplot[omThm,only marks,mark=*,mark size=1.6pt] table[x=dist,y=daily]{figdata/methods/02-brownian-motion/stops.csv};
\legend{continu~: réflexion,clôtures quotidiennes + bridge (sim.),quotidien~: niveau décalé,clôtures quotidiennes (sim.)}
\end{axis}
\end{tikzpicture}

\omcaption{Probabilité qu'un stop placé sous le cours d'entrée d'un titre dont la volatilité est de 30~\% soit touché
dans un délai de 21 jours de cotation. Surveillé en continu (trait plein) ou seulement aux clôtures (tirets,
avec le décalage de Broadie--Glasserman--Kou)~; les marques sont des simulations de 40\,000 mois, les marques vides
étant complétées par la probabilité de franchissement du pont entre deux clôtures. Données~: tutoriel du chapitre,
graine fixée.}\label{fig:qm:brownian-motion:stops}
\end{omfigure}

\section{Variation quadratique}

\begin{definition}[Variation quadratique]\label{def:qm:brownian-motion:qv}
La \emph{variation quadratique}\index{variation quadratique} d'un processus $X$ sur $[0, t]$ est
la limite en probabilité $[X]_t = \lim\sum_k (X_{t_{k+1}} - X_{t_k})^2$ sur les subdivisions $0 = t_0
< \dots < t_n = t$ dont le pas $\max_k(t_{k+1} - t_k)$ tend vers zéro, lorsqu'elle existe.
\end{definition}

\begin{theorem}[Variation quadratique du mouvement brownien]\label{thm:qm:brownian-motion:qv}
$[W]_t = t$~: les sommes convergent vers $t$ dans $L^2$, et presque sûrement le long des subdivisions dyadiques.
Par conséquent, la variation totale $\sum_k|W_{t_{k+1}} - W_{t_k}|$ tend vers l'infini.
\end{theorem}

\begin{proof}
Avec $\Delta_k = W_{t_{k+1}} - W_{t_k}$, $\E[\sum\Delta_k^2] = \sum(t_{k+1} - t_k) = t$, et par
l'indépendance $\Var(\sum\Delta_k^2) = \sum 2(t_{k+1} - t_k)^2 \le 2t\cdot\text{pas} \to 0$. Pour
les subdivisions dyadiques, les variances sont sommables et le lemme de Borel--Cantelli donne
la convergence presque sûre. Si la variation totale restait bornée par $V$ le long d'une suite, alors $\sum\Delta_k^2
\le V\max_k|\Delta_k| \to 0$ par continuité uniforme, d'où une contradiction.
\end{proof}

Pour une fonction régulière, la somme des carrés des accroissements s'annule~; pour une trajectoire brownienne, elle vaut le
temps écoulé, de façon déterministe, quelle que soit la trajectoire réalisée. Pour $\sigma W$, elle vaut $\sigma^2 t$, si bien qu'une
trajectoire observée de plus en plus finement révèle sa volatilité exactement, tandis que sa dérive reste
inobservable sur toute période fixée~: c'est la raison pour laquelle la volatilité réalisée (livre~1,
chapitre~25) s'estime bien mieux qu'un rendement espéré, et c'est aussi la raison pour laquelle le
bruit de microstructure du chapitre~21 joue un rôle aussi important. La \Cref{fig:qm:brownian-motion:qv}
échantillonne une trajectoire en $4$ à $65\,536$ points.

\begin{omfigure}
\begin{tikzpicture}
\begin{groupplot}[group style={group size=2 by 1,horizontal sep=1.8cm},
  width=7.4cm,height=4.8cm,xmode=log,log basis x=2,xlabel={points dans $[0,1]$},
  tick label style={font=\small},label style={font=\small},
  grid=major,grid style={gray!25},xmin=3,xmax=90000,
  table/col sep=comma]
\nextgroupplot[ylabel={$\sum(\Delta W)^2$},ymin=0.6,ymax=1.3]
\addplot[omDef,thick,mark=*,mark size=1.2pt] table[x=n,y=qv]{figdata/methods/02-brownian-motion/qv.csv};
\addplot[gray,dashed,domain=3:90000] {1};
\nextgroupplot[ylabel={$\sum|\Delta W|$},ymode=log,ymin=1,ymax=400]
\addplot[omThm,thick,mark=*,mark size=1.2pt] table[x=n,y=tv]{figdata/methods/02-brownian-motion/qv.csv};
\end{groupplot}
\end{tikzpicture}

\omcaption{Une trajectoire brownienne sur $[0,1]$ échantillonnée en $2^k$ points. À gauche~: la somme des carrés
des accroissements se stabilise autour de $t = 1$ (1.006 à 65\,536 points). À droite~: la somme des accroissements absolus
croît comme la racine carrée du nombre de points, sans limite. Données~: tutoriel du chapitre,
graine fixée.}\label{fig:qm:brownian-motion:qv}
\end{omfigure}

\section{Le pont brownien}\label{sec:qm:brownian-motion:bridge}

\begin{definition}[Pont brownien]\label{def:qm:brownian-motion:bridge}
Un \emph{pont brownien}\index{pont brownien} de $x$ à $y$ sur $[t_0, t_1]$ est un \omterm{def:qm:brownian-motion:bm}{mouvement brownien}
conditionné par $W_{t_0} = x$ et $W_{t_1} = y$~; de $0$ à $0$ sur $[0, T]$, il peut s'écrire
$W_t - (t/T)W_T$, un \omterm{def:qm:brownian-motion:bm}{processus gaussien} centré de covariance $s(T - t)/T$ pour $s \le t$.
\end{definition}

\begin{proposition}[Ce qui se passe entre deux observations]\label{prop:qm:brownian-motion:between}
Sachant $W_{t_0} = x$ et $W_{t_1} = y$, avec $h = t_1 - t_0$~: (i) $W_s$ pour $t_0 < s < t_1$ est
normale d'espérance $x + (s - t_0)(y - x)/h$ et de variance $(s - t_0)(t_1 - s)/h$, en particulier
$\mathcal N((x+y)/2, h/4)$ au point médian~; (ii) pour $\sigma W$ et un niveau $b$ inférieur à $x$
et à $y$,
\[
\P\bigl(\min_{t_0 \le s \le t_1}\sigma W_s \le b \bigm| \sigma W_{t_0} = x, \sigma W_{t_1} = y\bigr) =
\exp\Bigl(-\frac{2(x - b)(y - b)}{\sigma^2 h}\Bigr).
\]
\end{proposition}

\begin{proof}
(i) résulte du conditionnement gaussien du vecteur $(W_{t_0}, W_s, W_{t_1})$. Pour (ii), prendre $\sigma = 1$,
$t_0 = 0$, $x = 0$ et niveau $-c$ avec $c > 0$~; par le principe de réflexion $\P(\min \le -c, W_h
\in dy) = \P(W_h \in -2c - dy)$, donc la probabilité conditionnelle est le rapport des deux densités
normales, $\exp(-((2c + y)^2 - y^2)/2h) = \exp(-2c(c + y)/h)$, ce qui est la formule avec $x - b
= c$ et $y - b = c + y$.
\end{proof}

\begin{method}[Deux usages du pont en simulation]\label{met:qm:brownian-motion:bridge}
(i) \emph{Construction des trajectoires du grossier au fin}~: tirer $W_T = \sqrt T Z_0$, puis le point médian de chacun des
intervalles de la partie (i), en divisant les intervalles par deux jusqu'à ce que la grille soit assez fine. La trajectoire a la
bonne loi, et les premières normales fixent sa forme à grande échelle, ce que le chapitre~26 exploite avec
des points à faible discrépance. (ii) \emph{Surveillance d'une barrière entre les points de grille}~: simuler la trajectoire
aux seules dates d'observation, et pour chaque intervalle, multiplier les probabilités $1 - p_k$ de
\emph{ne pas} franchir, avec $p_k$ issue de la partie (ii). Le complément à 1 du produit est une estimation sans biais de
la probabilité de franchissement continu, sans raffinement de la grille.
\end{method}

Deux clôtures situées chacune à 1~\% au-dessus du stop, séparées d'un jour, avec une volatilité de 30~\%, laissent une probabilité de 57~\%
que le stop ait été touché dans l'intervalle~: $\exp(-2 \times 0.01^2/(0.09/252))$. La
simulation de la \cref{fig:qm:brownian-motion:stops} qui ajoute ces probabilités aux clôtures quotidiennes
retrouve la réponse continue, 81.6~\%, sans grille plus fine.

\section{Tutoriel~: stops, réflexions et variation quadratique}\label{tut:qm:brownian-motion:stops}

\begin{tutorial}
\textbf{Objectif.} Simuler des trajectoires browniennes de deux façons, vérifier le principe de réflexion et la
\omterm{def:qm:brownian-motion:qv}{variation quadratique}, et évaluer le stop-loss de l'accroche par surveillance continue et discrète.
\textbf{État final~:} les \cref{fig:qm:brownian-motion:paths,fig:qm:brownian-motion:reflect,fig:qm:brownian-motion:stops,fig:qm:brownian-motion:qv}
et les probabilités 81.6~\% et 72.0~\%.
\begin{tutsteps}
\item \textbf{Trajectoires}, par accroissements et par la construction en pont de la \cref{met:qm:brownian-motion:bridge}(i).
\omcode{code/firm/mcengine/firm_mcengine.py}{86}{101}{Trajectoires browniennes construites du grossier au fin par le pont.}\label{lst:qm:brownian-motion:bridge}
\item \textbf{Le stop}~: observer le logarithme du cours à un nombre donné de moments dans la journée, et
éventuellement ajouter la probabilité de franchissement du pont entre deux observations.
\omcode{code/methods/02-brownian-motion/python/qm_brownian.py}{37}{52}{Surveillance discrète d'un stop, avec et sans correction du pont.}\label{lst:qm:brownian-motion:touch}
\item \textbf{Exécuter} \texttt{stop\_problem()}, \texttt{convergence\_table()} et
\texttt{fig\_brownian.py}~; vérifier que les moteurs C++20 et Rust tirent les mêmes normales que
Python à partir de la même graine (\texttt{code/firm/mcengine/}).
\end{tutsteps}
\textbf{Pour aller plus loin.} Remplacer le stop fixe par un stop suiveur situé 2~\% sous le
maximum courant et déterminer sa probabilité de déclenchement~; affecter au logarithme du cours la dérive $-\sigma^2/2$ et
comparer la simulation avec la \cref{prop:qm:brownian-motion:driftpassage}.
\end{tutorial}

\section{Construction~: le moteur de Monte-Carlo, première étape}\label{bld:qm:brownian-motion:mcengine}

\begin{build}
\bfield{Rôle} Chaque simulation de la firme miniature tire ici ses trajectoires~: les trajectoires browniennes de ce chapitre,
les équations différentielles stochastiques du chapitre~4, la réduction de variance
et les points à faible discrépance du chapitre~26.
\bfield{Interface} \texttt{SplitMix64(seed)}, \texttt{NormalStream(seed)} (Box--Muller, cosinus
en premier)~; \texttt{brownian\_paths(n\_paths, n\_steps, T, seed, corr=None)}~;
\texttt{bridge\_paths(n\_paths, levels, T, seed)}~;
\texttt{bridge\_crossing\_probability(x0, x1, barrier, var)}~; en C++20 et Rust,
\texttt{brownian\_path}, \texttt{bridge\_path} et \texttt{bridge\_crossing\_probability} sur un
\texttt{NormalStream}.
\bfield{Règles} Le flux de référence (SplitMix64, Box--Muller) est identique dans les trois
langages~; les trajectoires Python en masse utilisent le PCG64 de NumPy~; les trajectoires corrélées passent par le facteur de Cholesky
de la matrice de corrélation~; les graines sont des arguments explicites, jamais un état global.
\bfield{Tests d'acceptation} \texttt{code/firm/mcengine/\{tests,cpp,rust\}}~: le vecteur de test
de SplitMix64~; les premières normales et une trajectoire à quatre pas égales dans les trois langages à $10^{-14}$~;
$\Var(W_T) = T$ et $\Cov(W_s, W_t) = \min(s,t)$ pour les deux constructions~; la formule de franchissement
par rapport à une simulation fine.
\bfield{Pour aller plus loin} Paires antithétiques~; un générateur à compteur dont les flux peuvent être répartis entre
fils d'exécution (chapitre~26).
\end{build}

\begin{omsources}
\item L.~Bachelier, «~Théorie de la spéculation~», \emph{Annales scientifiques de l'École
Normale Supérieure} 17, 1900.
\item N.~Wiener, «~Differential-space~», \emph{Journal of Mathematics and Physics} 2, 1923.
\item M.~D.~Donsker, «~An invariance principle for certain probability limit theorems~»,
\emph{Memoirs of the AMS} 6, 1951.
\item M.~Broadie, P.~Glasserman et S.~Kou, «~A continuity correction for discrete barrier
options~», \emph{Mathematical Finance} 7, 1997.
\item I.~Karatzas et S.~E.~Shreve, \emph{Brownian Motion and Stochastic Calculus},
Springer, 2\up{e} éd., 1991.
\end{omsources}

\section{Exercices}

\begin{exercise}[$\star$]\label{exo:qm:brownian-motion:1}
Calculer $\Cov(W_2, W_5)$, $\Var(W_5 - W_2)$ et la corrélation de $W_2$ et $W_5$.
\end{exercise}

\begin{exercise}[$\star$]\label{exo:qm:brownian-motion:2}
Calculer la probabilité qu'un \omterm{def:qm:brownian-motion:bm}{mouvement brownien} standard atteigne 1 avant le temps 1~?
\end{exercise}

\begin{exercise}[$\star$]\label{exo:qm:brownian-motion:3}
Une action a une volatilité annuelle de 20~\%. Calculer l'écart type de son rendement logarithmique
sur un jour et sur une semaine de cinq jours de cotation.
\end{exercise}

\begin{exercise}[$\star\star$]\label{exo:qm:brownian-motion:4}
Un stop se situe 5~\% sous le prix d'entrée d'une action de volatilité 30~\%. Calculer la probabilité qu'elle soit
touchée d'ici trois mois (63 jours de cotation) si on la surveille en continu, et si on la surveille
aux seuls cours de clôture~?
\end{exercise}

\begin{exercise}[$\star\star$]\label{exo:qm:brownian-motion:5}
Montrer que $W_t^2 - t$ est une \omterm{def:qm:probability-at-speed:martingale}{martingale}, puis en déduire le temps moyen que met $W$ pour quitter
l'intervalle $(-a, b)$, $a, b > 0$.
\end{exercise}

\begin{exercise}[$\star\star$]\label{exo:qm:brownian-motion:6}
Pour un \omterm{def:qm:brownian-motion:bridge}{pont brownien} de 0 à 0 sur $[0, 1]$, donner la loi de sa valeur en $t = \tfrac12$ et la
probabilité qu'il atteigne $0.5$.
\end{exercise}

\begin{exercise}[$\star\star\star$]\label{exo:qm:brownian-motion:7}
\emph{Programmation.} La somme des carrés des accroissements de $W$ sur $[0,1]$ en $n$ points a un écart type
égal à $\sqrt{2/n}$. Vérifier cela en simulant 2\,000 trajectoires pour $n = 4\,096$, et dire si les
fluctuations de la \cref{fig:qm:brownian-motion:qv} sont compatibles avec ce résultat.
\end{exercise}

\begin{exercise}[$\star\star\star$]\label{exo:qm:brownian-motion:8}
\emph{Repérer l'erreur.} «~Notre Monte-Carlo pour un contrat à extinction, qui s'annule si le cours
se négocie sous une barrière, simule les clôtures quotidiennes et les compare à la barrière~; le prix
a convergé à quatre chiffres.~» Corriger cela.
\end{exercise}

\section{Problème~: le stop-loss}

\begin{problem}[{Problème du week-end --- fréquence de déclenchement d'un stop à 2~\% dans le mois}]\label{pb:qm:brownian-motion:1}
Une action de volatilité annuelle de 30~\% est achetée avec un stop situé 2~\% sous le prix d'entrée. Le logarithme
de son cours est modélisé comme un \omterm{def:qm:brownian-motion:bm}{mouvement brownien} de volatilité $\sigma = 0.30$ et sans dérive, sur 21
jours de cotation sur une année de 252 jours.

\medskip\noindent\textbf{Partie~I --- Le stop continu.}
\begin{enumerate}
\item Quel est le niveau du stop $b$ en logarithme du cours, et l'écart type $\sigma\sqrt T$ sur le mois~?
\item Quelle est la probabilité que l'action termine le mois sous le stop~?
\item Quelle est la probabilité qu'elle touche le stop pendant le mois, surveillé en continu~?
\item Pourquoi le second nombre est-il exactement le double du premier~?
\item Avec la dérive $-\sigma^2/2$ qui fait du cours une \omterm{def:qm:probability-at-speed:martingale}{martingale}, que donne la \cref{prop:qm:brownian-motion:driftpassage}~?
\end{enumerate}

\medskip\noindent\textbf{Partie~II --- Surveillance aux clôtures.}
\begin{enumerate}[resume]
\item Que donne une simulation de 200\,000 mois de clôtures quotidiennes~?
\item Quel est le décalage de Broadie--Glasserman--Kou pour des clôtures quotidiennes, et quel est le stop continu équivalent~?
\item Quelle probabilité donne le niveau décalé~?
\item Pour une surveillance toutes les 15 minutes (26 fois par jour), que donnent la simulation et le décalage~?
\item Pourquoi une grille quatre fois plus fine ne réduit-elle que de moitié l'écart avec la réponse continue~?
\end{enumerate}

\medskip\noindent\textbf{Partie~III --- Le bridge.}
\begin{enumerate}[resume]
\item Deux clôtures consécutives sont chacune 1~\% au-dessus du stop. Quelle est la probabilité que le stop ait été touché entre elles~?
\item Que donne la simulation quotidienne lorsqu'on ajoute la probabilité de franchissement de chaque jour~?
\item Pourquoi cet estimateur est-il sans biais pour la probabilité continue~?
\item Quelle solution est la moins coûteuse pour une précision visée~: 64 observations par jour, ou des observations quotidiennes avec le bridge~?
\end{enumerate}

\medskip\noindent\textbf{Partie~IV --- Jugement.}
\begin{enumerate}[resume]
\item Avec une volatilité de 15~\%, quelle est la probabilité continue~?
\item Sur trois mois, quelles sont les probabilités continue et sur clôtures quotidiennes~?
\item Le stop est un ordre en attente sur la bourse~; quel chiffre s'applique, et qu'est-ce que le modèle ne prend pas en compte~?
\item Le desk vend une note structurée qui se désactive si la clôture quotidienne passe sous le niveau. Quel chiffre doit servir à la pricer~?
\item Énoncer le \emph{résultat nommé}~: la probabilité que le stop à 2~\% soit touché en un mois, en continu et sur les clôtures.
\item En une phrase~: pourquoi une barrière surveillée en continu a-t-elle bien plus de chances d'être touchée qu'une barrière surveillée chaque jour~?
\end{enumerate}
\end{problem}

\section{Questions d'entretien}

\begin{interviewq}[$\star$ \iqroles{trader, chercheur}]\label{iq:qm:brownian-motion:1}
Un \omterm{def:qm:brownian-motion:bm}{mouvement brownien} sans dérive part de 0. Quelle est la probabilité qu'il atteigne $-1$ avant $+2$~?
\end{interviewq}

\begin{interviewq}[$\star$ \iqroles{chercheur}]\label{iq:qm:brownian-motion:2}
Pourquoi la \omterm{def:qm:brownian-motion:qv}{variation quadratique} d'une trajectoire brownienne vaut-elle $t$, et pourquoi cela importe-t-il au trading~?
\end{interviewq}

\begin{interviewq}[$\star\star$ \iqroles{chercheur, trader}]\label{iq:qm:brownian-motion:3}
Quelle est l'espérance du maximum d'un \omterm{def:qm:brownian-motion:bm}{mouvement brownien} standard sur $[0,1]$~?
\end{interviewq}

\begin{interviewq}[$\star\star$ \iqroles{développeur, chercheur}]\label{iq:qm:brownian-motion:4}
Comment évaluez-vous une barrière surveillée en continu par Monte Carlo, sans recourir à de très petits pas de temps~?
\end{interviewq}

\begin{interviewq}[$\star\star$ \iqroles{risque, trader}]\label{iq:qm:brownian-motion:5}
Un trader affirme qu'un stop à 2~\% sur une action de volatilité de 30~\% est déclenché environ 40~\% du temps en un
mois. Qu'a-t-il calculé, et quel est le bon chiffre~?
\end{interviewq}

\begin{interviewq}[$\star\star\star$ \iqroles{chercheur}]\label{iq:qm:brownian-motion:6}
Calculez $\P(W_1 > 0, W_2 > 0)$.
\end{interviewq}
